\section{Begrensninger}
\label{sec:begrensninger}

VideoAPI er levert som et proof-of-concept. Det betyr at løsningen
demonstrerer at arkitekturen fungerer og at de definerte kravene kan
oppfylles, men den er ikke ferdig herdet for produksjon. Flere av valgene
underveis har vært bevisste kompromisser tilpasset prosjektets ramme, og
det er viktig å være tydelig på hva disse innebærer.

Den mest åpenbare begrensningen er tidsrammen. Med 19 uker tilgjengelig,
hvorav en betydelig andel gikk til å bygge grunnleggende forståelse av
videomotorene og WebRTC, ble det ikke rom for alt vi i utgangspunktet
hadde sett for oss. Validering av flere providere underveis kunne vært
utvidet med ekte tredjepartsmotorer i stedet for en mock-provider, og
ende-til-ende-testing kunne vært gjort mer systematisk. Vi måtte også
prioritere bredde framfor dybde på enkelte områder for å rekke å levere et
helhetlig system.

På sikkerhetssiden gjenstår det konkrete punkter som er beskrevet i
avsnitt~\ref{sec:eval-sikkerhet}.
Issuer- og audience-validering er deaktivert, \texttt{RequireHttpsMetadata}
er satt til \texttt{false}, og \texttt{/metrics}-endepunktet er eksponert
uten autentisering. Dette er bevisste valg knyttet til prosjektets
midlertidige infrastruktur og HDOs planlagte rekonfigurering, men det er
ikke valg man kan ta med seg uendret inn i produksjon. Autorisasjonsmodellen
er bevisst grovkornet, basert på et enkelt \texttt{api\_access}-claim, og må
utvides med rollestyring tilpasset HDOs egne behov før systemet brukes i
operativ sammenheng.

Avhengighetene til eksterne aktører setter også klare grenser for hva vi har
kunnet teste og verifisere. NHN-integrasjonen er bygget mot et QA-miljø, og
vi har opplevd både utdaterte credentials, lisensproblemer og URL-er i feil
format gjennom prosjektperioden. URL-rewrite-funksjonen i
\texttt{NHNVideoProvider} er en midlertidig løsning på et problem som egentlig
hører hjemme på NHN sin side. Når NHN retter formatet på sin ende vil
rewrite-koden bli overflødig, men inntil videre er den nødvendig for at
løsningen skal fungere. AntMedia-integrasjonen er testet mot HDOs
QA-instans, og avhenger av at konfigurasjonen og STUN/TURN-serverne der
fortsetter å være tilgjengelige.

Observability-oppsettet inneholder også flere pragmatiske kompromisser. UID
for Prometheus-datasource er hardkodet i Grafana-dashboardene,
\texttt{/metrics} unntas fra HTTPS-redirect for å la Prometheus scrape over
HTTP internt i Docker-nettverket, og portene for Grafana og Prometheus er
eksponert utad for at HDO skulle kunne teste dem direkte. Ingen av disse
valgene er gale i en utviklings- eller demonstrasjonskontekst, men de bør
vurderes på nytt før systemet settes i drift.

Testdekningen ligger på rundt 66 prosent linjedekning. Kjernelogikken og
provider-laget er godt dekket, men deler av infrastruktur- og
oppstartskoden er ikke testet direkte, og demo-frontenden er ikke en del av
testregimet. Vi har heller ikke gjort lasttesting eller målinger av reell
ytelse under produksjonslignende belastning. Det betyr at vi vet at API-et
fungerer korrekt for de scenariene vi har testet, men ikke hvor godt det
skalerer eller hvordan det oppfører seg under press.

Til slutt er det verdt å påpeke at demo-frontenden er nettopp en demo. Den
ble laget for å kunne vise at videostrømmer faktisk fungerte ende-til-ende,
ikke for å være et brukbart klientgrensesnitt. Den klienten HDO eventuelt
kobler på i produksjon vil måtte bygges separat, og en del av antagelsene i
demo-frontenden vil ikke nødvendigvis holde der.
